\documentclass[12pt,a4paper]{article}
\usepackage[utf8]{inputenc}
\usepackage[french]{babel}
\usepackage[T1]{fontenc}
\usepackage{amsmath}
\usepackage{amsfonts}
\usepackage{amssymb}
\usepackage{graphicx}
\usepackage{float}
\usepackage{algorithm}
\usepackage{algpseudocode}
\usepackage{tikz}
\usepackage{hyperref}
\usepackage{bm} % Pour les vecteurs en gras

% Configuration de la table des matières
\usepackage{tocloft}
\renewcommand{\cftsecleader}{\cftdotfill{\cftdotsep}}
\renewcommand{\cftsecfont}{\normalfont\bfseries}
\renewcommand{\cftsubsecfont}{\normalfont}

% Configuration des marges
\usepackage[top=2.5cm,bottom=2.5cm,left=2.5cm,right=2.5cm]{geometry}

% Configuration des liens hypertexte
\hypersetup{
    colorlinks=true,
    linkcolor=blue,
    citecolor=red,
    urlcolor=cyan,
    pdftitle={Roguia: Documentation Technique IA},
    pdfauthor={Messaoudi, Ringuet, Ravirooban, Mbaye, Klein}
}

% Commandes mathématiques personnalisées
\newcommand{\state}{\mathbf{s}}
\newcommand{\action}{a}
\newcommand{\reward}{r}
\newcommand{\policy}{\pi}
\newcommand{\Qval}{Q}
\DeclareMathOperator*{\argmax}{arg\,max}

\title{
    \vspace{2cm}
    \textbf{\Huge Roguia : Une Approche par Apprentissage par Renforcement Hybride dans un Roguelike} \\ [1cm]
    \large \textit{Documentation Technique et Formalisation Mathématique}
}
\author{
    \textbf{Équipe Projet Roguia} \\
    ESEO - Option DAIA \\
    \textit{Supervisé par A. Petrossian \& R. Costaouec}
}
\date{\today}

\begin{document}

\maketitle
\thispagestyle{empty}
\newpage

\tableofcontents
\newpage

\section{Introduction}
Ce document présente la formalisation mathématique et l'implémentation technique des agents autonomes au sein du projet \textit{Roguia}. L'objectif est de développer une intelligence artificielle capable de contrôler des entités ennemies ("monstres") dans un environnement stochastique généré procéduralement.

L'approche retenue repose sur l'Apprentissage par Renforcement (RL), et plus spécifiquement sur l'algorithme Q-Learning \cite{watkins1989learning}. Le système se distingue par une architecture hybride combinant un apprentissage en ligne (\textit{Online Learning}) local à chaque instance de jeu, et un apprentissage hors ligne (\textit{Offline/Batch RL}) globalisé, permettant l'émergence d'un "Super Monstre" via l'agrégation de données massives stockées sur MongoDB.

\end{document}